\chapter{Conclusion}
\label{ch:conclusion}

\section{Summary of Findings}

This project designed, implemented, and deployed CAAS --- a complete end-to-end
MLOps pipeline for PM2.5 air quality forecasting in Chiang Mai, Thailand.

The principal findings are:

\begin{enumerate}
    \item \textbf{XGBoost outperforms LSTM at all three horizons.} The champion
          XGBoost model achieves MAE of 5.31, 6.91, and 8.34~\textmu g/m\textsuperscript{3}
          at t+1, t+3, and t+7 respectively. It is interpretable via SHAP, fast
          to train, and more stable than the LSTM across horizons. The LSTM
          provides competitive results at t+7 (MAE 8.06 vs.\ 8.34) but degrades
          significantly at t+3 (R²=0.54 vs.\ 0.74), confirming that tabular
          XGBoost with engineered features is the superior approach for this problem.

    \item \textbf{NASA FIRMS fire hotspot data adds measurable value.} Removing
          the six fire-derived features increases MAE by 4.3--6.8\% across horizons.
          The effect is strongest at t+1, where recent hotspot counts directly
          predict next-day PM2.5 accumulation. This validates the cost of the
          FIRMS data integration.

    \item \textbf{The predictive signal shifts with horizon.} SHAP analysis reveals
          a clear transition: at t+1, recent PM2.5 (\texttt{pm25\_lag1}) dominates;
          at t+3, 14-day fire accumulation (\texttt{hotspot\_14d\_roll}) is most
          important; at t+7, seasonal position (\texttt{sin\_month}) is the
          strongest signal. This is physically interpretable and consistent with
          the meteorology of Chiang Mai haze events.

    \item \textbf{The default 50~\textmu g/m\textsuperscript{3} alert threshold is
          near-optimal for multi-day horizons.} Scenario~C shows that only the
          t+1 classifier benefits meaningfully from threshold tuning
          (53.7~\textmu g/m\textsuperscript{3}, +4.0\% F1). The standard WHO/Thai
          threshold is well-calibrated for t+3 and t+7.

    \item \textbf{The MLOps pipeline is production-ready within its design scope.}
          Automated daily ingestion, drift monitoring with a seasonal-aware retrain
          policy, champion/challenger validation, and Terraform-provisioned AWS
          infrastructure together constitute a deployable, maintainable system.
\end{enumerate}

\section{Limitations}

\begin{itemize}
    \item \textbf{Single station.} CAAS currently operates on one Chiang Mai
          monitoring station. PM2.5 varies spatially, and a network of stations
          would provide better coverage for the population.
    \item \textbf{Historical weather only.} The pipeline uses ERA5 reanalysis
          weather data (accurate but historical). True production use requires
          weather forecast inputs (e.g.\ GFS or ECMWF), which introduce their own
          forecast uncertainty.
    \item \textbf{LSTM underperformance at t+3.} Despite architectural improvements
          (seq\_len=60, deeper network, BatchNorm), the LSTM R² of 0.54 at t+3
          indicates the model struggles at the mid-range horizon. This suggests
          the temporal sequences do not provide sufficient additional signal over
          the tabular lag features already engineered for XGBoost.
    \item \textbf{Sparse fire data in non-haze months.} Zero-inflation in the FIRMS
          features during the rainy season limits PSI-based drift detection and
          requires suppression logic. A seasonal model could handle this more
          elegantly.
\end{itemize}

\section{Future Work}

\begin{itemize}
    \item \textbf{Multi-station expansion:} Extend to all 60+ PCD monitoring
          stations in Thailand, enabling national-scale air quality forecasting.
    \item \textbf{Weather forecast integration:} Replace reanalysis inputs with
          ECMWF HRES forecast data to enable genuine real-time t+3 and t+7
          predictions without relying on observed weather.
    \item \textbf{Graph Neural Network:} Model spatial dependencies between
          nearby stations using a GNN, potentially improving long-horizon
          forecasts by propagating upstream fire signals across the network.
    \item \textbf{Hyperparameter optimisation:} Replace manual tuning with Optuna
          or Ray Tune to systematically search the hyperparameter space.
    \item \textbf{Kubernetes serving:} For production deployment at scale, migrate
          from a single EC2 instance to a Kubernetes cluster with horizontal
          pod autoscaling.
\end{itemize}
